\chapter{Discussion}
\label{ch:discussion}

This chapter interprets the results at the level of pattern and mechanism --- the parts that
are already understood from the design and the validation runs --- and states the threats to
validity honestly. The quantitative confirmations are deferred to the filled tables of
\cref{ch:results}; the qualitative claims here are stable.

\section{What the evaluation is really measuring}

The thesis's contribution is not a profitable agent but a \emph{measurement}. The combination
of near-perfect grounding with a regime-dependent faithfulness gap is the central qualitative
result, and it is more informative than either metric alone. High grounding establishes that
the agent is not hallucinating its inputs: when it cites a moving average or a position
weight, that number is real. A faithfulness gap on top of high grounding therefore cannot be
explained away as ``the model made up the numbers''; it is a genuine knowledge--action gap of
the kind KellyBench~\cite{kellybench2026} predicts --- the agent sees the right evidence,
cites it correctly, and then acts against it.

\section{The exit-discipline gap}

The dominant mechanism behind unfaithful decisions is a failure to trim positions that have
grown past the cap, concentrated in the strongest-trending names and in the bull regime
(\cref{tab:nontrim}). This is a specific, interpretable behaviour rather than a diffuse failure,
and its diagnosis is itself a methodological result. It emerged only because the faithfulness
judge was made policy-aware: the naive judge silently reclassified these decisions as
``constrained,'' hiding the gap entirely. The lesson generalises beyond this project: a metric
that bakes in the wrong assumption about the agent's policy can mask exactly the behaviour it
is meant to surface, and the discipline of \emph{pre-registering} the policy-aware metric
before computing it is what makes the resulting finding credible rather than a post-hoc story.

The behaviour also exposed a prompt ambiguity worth naming. The original sizing instructions
constrained entry but were less explicit about trimming an appreciated position; the imperative
that a position ``must not exceed'' the cap could be read either as an ongoing maintenance rule
or as an entry-only constraint. The common-reference run resolves that ambiguity operationally:
the system prompt now states that positions above $20\%$ \nav\ must be trimmed. This makes any
remaining non-trim gap easier to interpret as behaviour rather than instruction ambiguity.

\section{Single-agent versus multi-agent}

Following AMA~\cite{ama2025}, the multi-agent Portfolio Manager is evaluated for behavioural
change rather than assumed superiority. The common-reference run shows the conservative
outcome: explicit deliberation generalises the causal, grounded evaluation harness to a richer
architecture, but it does not automatically produce a statistically reliable performance edge
or close every action-consistency gap. This strengthens the thesis's evaluation-first stance:
more architecture is not automatically more capability, and the harness is valuable precisely
because it can show that clearly.

\section{Limitations and threats to validity}
\label{sec:limitations}

We state limitations plainly; for a thesis built on rigour, candour about what the design does
\emph{not} establish is part of the contribution.

\paragraph{Weight-memory contamination.} The most serious threat is unremovable by any data
filter: the model's parameters encode the approximate outcome of the historical test period
(\cref{sec:leakage}). The \tnow\ filter prevents \emph{data} leakage but not \emph{parametric}
recall. We mitigate by instructing a rule-based process over tools and, where feasible,
pushing the window toward the training cutoff, but we do not claim to eliminate it. A real news
corpus does not solve this; it only makes grounding meaningful. This caps the strength of any
claim that the agent's reasoning generalises to genuinely unseen periods.

\paragraph{Single market window and asset class.} The primary study is US equities over one
multi-regime window. The regime labels are causal and the segmentation is principled, but four
regime buckets in one window are not a substitute for multiple independent periods. Cross-asset
and out-of-period generalisation are future work, not claims of this thesis.

\paragraph{Bootstrap assumptions.} The confidence intervals use an i.i.d.\ bootstrap; the
stationary bootstrap of Politis and Romano~\cite{politis1994} would better respect return
autocorrelation, and our intervals may be mildly optimistic for momentum-laden strategies
(\cref{sec:bootstrap}). We report this rather than hide it.

\paragraph{LLM-judge dependence of faithfulness.} Faithfulness rests on a judge that is itself
an LLM. We mitigate with a blind protocol, temperature-zero determinism, pre-registration, and
human anchoring (\cref{sec:anchor}), but the metric is ultimately a model-mediated estimate of
a counterfactual, and the residual judge weakness around the entry threshold is documented.

\paragraph{Unadjusted prices.} Avoiding split-adjustment leakage costs us split-corrected
long-horizon comparisons (\cref{sec:leakage}). This is the right trade-off for a look-ahead
study but does slightly distort multi-year return arithmetic.

\paragraph{Annualisation and short windows.} Validation-window performance is never
annualised, by design; any short-window Sharpe is statistically meaningless and is excluded
from all claims.

\section{Future work}
\label{sec:future}

Three extensions are identified and deliberately left out of scope to protect the evaluation.
First, a \emph{quant-researcher} agent that reads the literature and registers custom
indicators into the analytics server, extending the agent's own toolbox --- an original angle
relative to TradingAgents, but a scope risk that must be sandboxed and validated before any
generated code is executed. Second, \emph{reinforcement-learning fine-tuning} using the
decision and realised-return logs as a reward signal, contingent on a strictly held-out
out-of-sample period to detect overfitting. Third, the methodological upgrades already flagged:
an HMM regime model (causally fitted), the stationary bootstrap, and a cross-provider backbone
comparison in the spirit of~\cite{ama2025}. Each is valuable; none is a prerequisite for the
thesis's contribution, and each is recorded so that a successor can pick it up cleanly.
